\chapter{范畴：对象、态射、复合、函子}\label{ux8303ux7574ux5bf9ux8c61ux6001ux5c04ux590dux5408ux51fdux5b50}

前四章处理 Haskell 的语法与函数。本章引入范畴的定义，并说明它刻画的正是第 2 章已经用到的
复合运算。

\section{从复合律出发}\label{ux4eceux590dux5408ux5f8bux51faux53d1}

第 2 章的 \passthrough{\lstinline!(.)!} 与 \passthrough{\lstinline!id!} 满足三条等式：

\begin{lstlisting}[language=Haskell]
id . f        = f            -- 左单位
f . id        = f            -- 右单位
(f . g) . h   = f . (g . h)  -- 结合律
\end{lstlisting}

\passthrough{\lstinline!Control.Category!} 把这些等式列为定律，但不作检查：实例写错不会报错，只会产生难以定位的
结果。\passthrough{\lstinline!fmap!}、单子变换等大量函数都建立在这三条等式之上。

因此需要确定：满足这三条等式的结构在数学上叫什么。答案是\textbf{范畴}（category）。

\section{范畴的定义}\label{ux8303ux7574ux7684ux5b9aux4e49}

\textbf{定义 5.1（范畴）} 一个范畴 \(\mathcal{C}\) 由三部分组成：

\begin{enumerate}
\def\labelenumi{\arabic{enumi}.}
\tightlist
\item
  \textbf{对象}的汇集 \(\mathrm{Ob}(\mathcal{C})\)；
\item
  对每一对对象 \(A, B\)，一个\textbf{态射}的汇集 \(\mathrm{Hom}_{\mathcal{C}}(A, B)\)，
  其中的元素写作 \(f \colon A \to B\)；
\item
  对每个对象 \(A\)，一个\textbf{恒等态射} \(1_A \colon A \to A\)；对每一对可复合的态射
  \(f \colon A \to B\)、\(g \colon B \to C\)，一个\textbf{复合} \(g \circ f \colon A \to C\)。
\end{enumerate}

并要求两条公理成立：

\begin{itemize}
\tightlist
\item
  \textbf{结合律}：对 \(f \colon A \to B\)、\(g \colon B \to C\)、\(h \colon C \to D\)，
  \[h \circ (g \circ f) = (h \circ g) \circ f\]
\item
  \textbf{单位律}：对 \(f \colon A \to B\)，
  \[1_B \circ f = f = f \circ 1_A\]
\end{itemize}

定义不涉及元素、加法或拓扑。它刻画的是一个带类型的箭头系统，以及该系统上一个可结合的复合
运算。

结构如图：

\begin{lstlisting}
        f          g
   A ------> B ------> C
   |          |          |
  1_A       1_B        1_C
   |          |          |
   v          v          v
   A          B          C

   g . f : A → C 表示先 f 后 g
\end{lstlisting}

两条公理的作用：结合律保证 \passthrough{\lstinline!f . g . h!} 的括号位置不影响结果，单位律保证每个对象上存在一个
不改变结果的中性箭头。缺少任何一条，复合都无法作为基本操作使用。

🦂 数学写作 \(g \circ f\)，读作“先 \(f\) 后 \(g\)”，即 \((g \circ f)(x) = g(f(x))\)。Haskell 的
\passthrough{\lstinline!g . f!} 与它是同一个东西：\passthrough{\lstinline!(g . f) x = g (f x)!}，不需要换算。

容易读错的地方不在两种记号之间，而在记号与执行顺序之间：复合写在左边的函数后执行。
\passthrough{\lstinline!f . g . h!} 从右往左读，依次执行 \passthrough{\lstinline!h!}、\passthrough{\lstinline!g!}、\passthrough{\lstinline!f!}，与 \(f \circ g \circ h\) 相同。

\section{例子}\label{ux4f8bux5b50}

\subsection{例一：Hask}\label{ux4f8bux4e00hask}

对象是 Haskell 的类型，态射是函数，复合是 \passthrough{\lstinline!(.)!}，恒等是 \passthrough{\lstinline!id!}。

\begin{lstlisting}[language=Haskell]
Int      -- 对象
not      -- 态射 Bool -> Bool
(.)      -- 复合
id       -- 恒等
\end{lstlisting}

这个范畴通常记作 \textbf{Hask}。它在 \(\bot\) 与 \passthrough{\lstinline!seq!} 介入时并不严格满足公理，第 6 章处理这一
问题；本章先按满足处理。

\subsection{例二：偏序集}\label{ux4f8bux4e8cux504fux5e8fux96c6}

设 \((P, \leq)\) 为偏序集，构造范畴如下：

\begin{itemize}
\tightlist
\item
  对象：\(P\) 中的元素；
\item
  态射：当 \(a \leq b\) 时，\(\mathrm{Hom}(a, b)\) 中恰有一个箭头 \(a \to b\)；否则为空。
\end{itemize}

恒等箭头对应自反性 \(a \leq a\)；复合对应传递性：由 \(a \leq b\) 与 \(b \leq c\) 得 \(a \leq c\)。
结合律自动成立，因为同一对对象之间至多一个箭头，两条路径必然给出同一个箭头。

这类范畴的对象多而箭头少：\(n\) 个元素至多 \(n + \binom{n}{2}\) 个箭头，且每对对象之间至多一个。
信息集中在箭头是否存在。

\subsection{例三：幺半群}\label{ux4f8bux4e09ux5e7aux534aux7fa4}

设 \((M, \cdot, e)\) 为幺半群，构造范畴如下：

\begin{itemize}
\tightlist
\item
  对象：只有一个，记作 \(\ast\)；
\item
  态射：\(M\) 中每个元素 \(m\) 对应一个箭头 \(\ast \to \ast\)；
\item
  复合：\(m \circ n = m \cdot n\)；
\item
  恒等：\(e\)。
\end{itemize}

这类范畴的对象只有一个，信息全部集中在箭头的复合方式上。

公理有实际约束力。浮点加法不满足结合律：

\begin{lstlisting}[language=Haskell]
(0.1 + 0.2) + 0.3            -- 0.6000000000000001
0.1 + (0.2 + 0.3)            -- 0.6
\end{lstlisting}

两者不等，因此 \((\mathrm{Double}, +, 0)\) 不是幺半群，也不能按上述方式构成单对象范畴。
这一差别有实际后果：若复合运算是浮点累加，任何依赖括号重排的变换（例如并行求和）都可能改变
结果。

\subsection{例四：集合范畴}\label{ux4f8bux56dbux96c6ux5408ux8303ux7574}

所有集合与它们之间的全部函数构成一个范畴，记作 \textbf{Set}：

\begin{itemize}
\tightlist
\item
  对象：集合；
\item
  态射：集合之间的函数；
\item
  复合：函数的复合；
\item
  恒等：恒等函数。
\end{itemize}

Set 是数学中最常用的范畴。唯一的规模问题：全体集合的汇集不是集合，而是\textbf{真类（proper
class）}，因此 Set 的“对象全体”比任何集合都大。本书不在规模问题上停留，只在 6.1 节与
Hask 对照时用到这一区别。

\subsection{两类极端与 Set}\label{ux4e24ux7c7bux6781ux7aefux4e0e-set}

偏序集与幺半群是范畴的两个极端：前者把结构放在对象之间（箭头是否存在），后者把结构放在
箭头之间（如何复合）。Set 位于两者之间：对象多，箭头也多，任意两个集合之间的函数仍构成一个集合。一般范畴同时具有两个方向，因此同一概念在这三类范畴中呈现的形态不同。

这一对比有助于理解后文：

\begin{itemize}
\tightlist
\item
  积在偏序集中是最大下界（当它存在时）；在单对象范畴中一般不存在：唯一对象只能与自身配对，
  而积要求任意一对元素 \((f, g)\) 都唯一地写成 \((\pi_1 h, \pi_2 h)\)，即乘法给出
  \(M \times M\) 到 \(M\) 的双射，有限幺半群只有平凡情形做得到。积的一般定义见第 7 章；
\item
  单子不是单对象范畴中的东西：它是自函子范畴中的幺半群，即一个自函子加两个自然变换，见第 9
  章关于自函子范畴的讨论。
\end{itemize}

五个例子的对照见表 5.1：

{\small\setlength{\tabcolsep}{4pt}\def\LTcaptype{table} % do not increment counter
\begin{longtable}[]{@{}
  l
  l
  l
  l
  l@{}}
\caption{五个例子的对照}\label{tab:5-1}\\
\toprule\noalign{}
\textbf{范畴} & \textbf{对象} & \textbf{态射} & \textbf{复合} & \textbf{恒等} \\
\midrule\noalign{}
\endhead
\bottomrule\noalign{}
\endlastfoot
Hask & Haskell 类型 & 函数 & \passthrough{\lstinline!(.)!} & \passthrough{\lstinline!id!} \\
偏序集 \((P, \leq)\) & \(P\) 的元素 & \(a \le b\) 时恰好一个 & 传递性 & 自反性 \\
幺半群 \((M, \cdot, e)\) & 唯一对象 \(\ast\) & \(M\) 的元素 & 乘法 & \(e\) \\
Set & 集合 & 函数 & 函数复合 & 恒等函数 \\
三对象范畴 & \(A, B, C\) & 六个箭头 & 查表 & \(1_A, 1_B, 1_C\) \\
\end{longtable}
}

同一个定义覆盖五类不同的结构。范畴论抽象出的是复合运算本身，而不是被复合的内容。

\subsection{同构}\label{ux540cux6784}

有了复合运算，可以定义同构，它是范畴中“两对象相同”的精确含义。

\textbf{定义 5.2（同构）} 态射 \(f \colon A \to B\) 称为\textbf{同构}，若存在 \(g \colon B \to A\) 使得

\[g \circ f = 1_A, \qquad f \circ g = 1_B\]

此时记 \(A \cong B\)。在 Hask 中，同构即两个函数互为反函数。三个例子：

\begin{itemize}
\tightlist
\item
  \passthrough{\lstinline!swapPair :: (a, b) -> (b, a)!} 自逆，做两次回到原值；
\item
  \passthrough{\lstinline!curry!} 与 \passthrough{\lstinline!uncurry!} 互为反函数，即 2.4 节讨论的双射；
\item
  \passthrough{\lstinline!newtype Age = Age Int!} 的构造与拆解。
\end{itemize}

\begin{lstlisting}[language=Haskell]
-- | 交换二元组：做了两次就回到原样，所以它是自逆的。
swapPair :: (a, b) -> (b, a)
swapPair (x, y) = (y, x)

-- | 同构的判据：@swapPair . swapPair = id@。
swapIsIso :: Bool
swapIsIso = and
    [ swapPair (swapPair p) == p
    | p <- [(1 :: Int, 'a'), (2, 'b'), (3, 'c')]
    ]
\end{lstlisting}

\passthrough{\lstinline!swapIsIso!} 在测试中为 \passthrough{\lstinline!True!}。

同构弱于相等：\(A\) 与 \(B\) 不必是同一个对象，只需在范畴内部无法区分。这一点在 Hask 中很具体：
\passthrough{\lstinline!(a, b)!} 与 \passthrough{\lstinline!(b, a)!} 不是同一个类型，但作为范畴对象同构，因此任何用函数与复合表达的性质必然
相同。

\section{用有限范畴验证公理}\label{ux7528ux6709ux9650ux8303ux7574ux9a8cux8bc1ux516cux7406}

下面构造一个三对象范畴，并逐条验证公理。

三个对象 \(A, B, C\)，箭头构成一条链加一条捷径：

\begin{center}
\begin{tikzpicture}[node distance=3cm, auto]
  \node (A) at (0,0) {$A$};
  \node (B) at (3,0) {$B$};
  \node (C) at (6,0) {$C$};
  \draw[->] (A) -- node[above] {$F$} (B);
  \draw[->] (B) -- node[above] {$G$} (C);
  \draw[->] (A) to [bend right=35] node[below] {$H$} (C);
\end{tikzpicture}
\end{center}

再加三个恒等箭头，共六个。对象定义为枚举：

\begin{lstlisting}[language=Haskell]
data Obj = A | B | C
    deriving (Eq, Show, Enum, Bounded)
\end{lstlisting}

箭头同样是枚举，包含三个恒等与三个非恒等箭头：

\begin{lstlisting}[language=Haskell]
-- | Six arrows: three identities, three "real" ones.
data Arrow = IdA | IdB | IdC | F | G | H
    deriving (Eq, Show, Enum, Bounded)

allArrows :: [Arrow]
allArrows = [minBound .. maxBound]
\end{lstlisting}

恒等由对象确定：

\begin{lstlisting}[language=Haskell]
identity :: Obj -> Arrow
identity A = IdA
identity B = IdB
identity C = IdC
\end{lstlisting}

复合通过查表实现。注意 \passthrough{\lstinline!compose g f!} 的第一个参数是后执行的态射，即 \(g \circ f\)：

\begin{lstlisting}[language=Haskell]
-- | @compose g f@ is the arrow "first @f@, then @g@", written @g . f@ in maths.
-- Returns 'Nothing' when the types do not line up.
compose :: Arrow -> Arrow -> Maybe Arrow
compose g f = lookup (g, f) table

table :: [((Arrow, Arrow), Arrow)]
table =
    [ ((IdA, IdA), IdA)
    , ((F,   IdA), F)
    , ((H,   IdA), H)
    , ((IdB, IdB), IdB)
    , ((G,   IdB), G)
    , ((IdC, IdC), IdC)
    , ((IdB, F),   F)
    , ((G,   F),   H)
    , ((IdC, G),   G)
    , ((IdC, H),   H)
    ]
\end{lstlisting}

表只有十项，因为只有十对态射可以复合。\passthrough{\lstinline!compose G F!} 为 \passthrough{\lstinline!Just H!}，即先 \(F\) 后 \(G\)，正是捷径；
\passthrough{\lstinline!compose F G!} 为 \passthrough{\lstinline!Nothing!}：\(F\) 的定义域是 \(A\)，而 \(G\) 的值域是 \(C\)，没有 \(C \to A\) 的箭头。
这与类型检查的作用相同，只是在此手写。

公理可以写成可执行的断言：

\begin{lstlisting}[language=Haskell]
-- | Associativity, checked on every triple: @(h . g) . f = h . (g . f)@.
assocHolds :: Bool
assocHolds = and
    [ composeM (compose h g) (Just f) == composeM (Just h) (compose g f)
    | h <- allArrows, g <- allArrows, f <- allArrows
    ]

-- | @1 . f = f@ and @f . 1 = f@ for every arrow @f@.
identityHolds :: Bool
identityHolds = and
    [ compose (identity (tgt f)) f == Just f
      && compose f (identity (src f)) == Just f
    | f <- allArrows
    ]
\end{lstlisting}

\passthrough{\lstinline!assocHolds!} 枚举全部 \(6^3 = 216\) 个三元组，\passthrough{\lstinline!identityHolds!} 检查六个箭头的左右单位律。
两者在测试中均为 \passthrough{\lstinline!True!}。

同一模式可用于幺半群：

\begin{lstlisting}[language=Haskell]
-- | The monoid @Z/3@ seen as a one-object category: one object, three arrows,
-- composition is addition modulo 3, the identity is zero.
data M3 = M0 | M1 | M2
    deriving (Eq, Show, Enum, Bounded)

madd :: M3 -> M3 -> M3
madd x y = toEnum ((fromEnum x + fromEnum y) `mod` 3)

assocHoldsM3 :: Bool
assocHoldsM3 = and
    [ madd (madd x y) z == madd x (madd y z)
    | x <- [M0 .. M2], y <- [M0 .. M2], z <- [M0 .. M2]
    ]
\end{lstlisting}

\(Z/3\) 在模 3 加法下满足结合律。这原是一个算术事实，现在也是“某个单对象范畴满足结合律”的
实例。同一段检查代码适用于两个来源不同的范畴，原因是范畴论刻画的是复合运算，而非复合的
具体内容。

◇ \textbf{恒等箭头唯一。} 若 \(e, e'\) 都是 \(A\) 上的恒等，则 \(e = e \circ e' = e'\)：第一处用
\(e'\) 的恒等性质，第二处用 \(e\) 的恒等性质。因此定义中说的是“一个恒等”而非“至少一个”。

\section{函子}\label{ux51fdux5b50}

有了范畴，接下来是范畴之间的比较，对应的概念是\textbf{函子}（functor）。

\textbf{定义 5.3（函子）} 从 \(\mathcal{C}\) 到 \(\mathcal{D}\) 的函子 \(F\) 由两部分组成：

\begin{itemize}
\tightlist
\item
  对象上的映射：\(A \mapsto FA\)；
\item
  态射上的映射：\(f \colon A \to B \mapsto Ff \colon FA \to FB\)。
\end{itemize}

并要求保恒等与保复合：

\[F(1_A) = 1_{FA}, \qquad F(g \circ f) = Fg \circ Ff\]

第二条的方向需要留意：左边是先把复合映射过去，右边是映射过去之后再复合。两者相等，即函子不破坏
结构的准确含义。

以上面的三对象范畴为例，将它映入 Hask：三个对象都映到 \passthrough{\lstinline!Int!}，三个箭头映到三个函数：

\begin{lstlisting}[language=Haskell]
-- | A functor from our three-object category into Hask: all three objects go
-- to 'Int', and the arrows go to functions.
interp :: Arrow -> Int -> Int
interp IdA = id
interp IdB = id
interp IdC = id
interp F   = (+ 1)
interp G   = (* 2)
interp H   = (* 2) . (+ 1)

-- | Functor laws: identity is preserved, composition is preserved.
functorLawHolds :: Bool
functorLawHolds = and
    [ identityOK x && compositionOK x g f
    | g <- allArrows, f <- allArrows, x <- [-2 .. 2 :: Int]
    ]
  where
    identityOK x = and [interp (identity o) x == x | o <- [A, B, C]]
    compositionOK x g f = case compose g f of
        Nothing -> True
        Just h  -> interp h x == interp g (interp f x)
\end{lstlisting}

由于 \(H = G \circ F\)，\(FH\) 必须等于 \(FG \circ FF\)，即 \passthrough{\lstinline!(*2) . (+1)!}。测试 \passthrough{\lstinline!functorLawHolds!}
在五个输入上逐条验证：

\begin{lstlisting}[language=Haskell]
interp H 3                 -- 8
interp G (interp F 3)      -- 也是 8
\end{lstlisting}

\subsection{一个反例}\label{ux4e00ux4e2aux53cdux4f8b}

将同一映射改动一处：\(H\) 不映到 \passthrough{\lstinline!(*2) . (+1)!}，而映到 \passthrough{\lstinline!(+3)!}：

\begin{lstlisting}[language=Haskell]
-- | 一个"看起来像函子"的映射：只把 @H@ 送错了一个函数。
badInterp :: Arrow -> Int -> Int
badInterp IdA = id
badInterp IdB = id
badInterp IdC = id
badInterp F   = (+ 1)
badInterp G   = (* 2)
badInterp H   = (+ 3)   -- 错的：H = G . F，应该等于 (* 2) . (+ 1)

-- | 同一套检查，作用在假函子上：返回 'False'。
badFunctorLawHolds :: Bool
badFunctorLawHolds = and
    [ case compose g f of
        Nothing -> True
        Just h  -> badInterp h x == badInterp g (badInterp f x)
    | g <- allArrows, f <- allArrows, x <- [-2 .. 2 :: Int]
    ]
\end{lstlisting}

它满足表面要求：三个对象都映到 \passthrough{\lstinline!Int!}，三个箭头都是 \passthrough{\lstinline!Int -> Int!}，恒等也保持。但
\passthrough{\lstinline!badFunctorLawHolds!} 为 \passthrough{\lstinline!False!}：取 \(x = 0\)，\passthrough{\lstinline!badInterp H 0!} 得 \passthrough{\lstinline!3!}，
\passthrough{\lstinline!badInterp G (badInterp F 0)!} 得 \passthrough{\lstinline!2!}。

函子定律的作用即在于此：区分“像函子”与“是函子”。Haskell 中 \passthrough{\lstinline!fmap!} 的实例写错时（例如把
\passthrough{\lstinline!fmap f (Just x)!} 写成 \passthrough{\lstinline!Just x!}），编译器不会报错，只有定律能查出。

第 7 章将说明，\passthrough{\lstinline!fmap!} 正是某个函子在态射上的部分。届时 \passthrough{\lstinline!fmap id = id!} 与
\passthrough{\lstinline!fmap (f . g) = fmap f . fmap g!} 无须单独记忆，它们是定义 5.3 的直接推论。

\section{常见错误}\label{ux5e38ux89c1ux9519ux8bef}

\subsection{复合方向读反}\label{ux590dux5408ux65b9ux5411ux8bfbux53cd}

\passthrough{\lstinline!f . g!} 表示先 \passthrough{\lstinline!g!} 后 \passthrough{\lstinline!f!}（见 5.2 节）。读反的原因不在 Haskell 的记号，而在书写顺序：
复合写在左边的函数后执行，与 \passthrough{\lstinline!f(x)!} 中“函数在左、先作用”的习惯相反。

可靠的检查方式是看类型：\passthrough{\lstinline!f . g!} 中 \passthrough{\lstinline!g!} 的输出必须作为 \passthrough{\lstinline!f!} 的输入，因此 \passthrough{\lstinline!g :: a -> b!}、
\passthrough{\lstinline!f :: b -> c!}。类型签名是唯一可靠的方向依据。

\subsection{把类型类当作对象}\label{ux628aux7c7bux578bux7c7bux5f53ux4f5cux5bf9ux8c61}

\passthrough{\lstinline!Eq!}、\passthrough{\lstinline!Show!}、\passthrough{\lstinline!Monad!} 是类型类，不是范畴的对象。在 Hask 中：

\begin{itemize}
\tightlist
\item
  对象是类型，如 \passthrough{\lstinline!Int!}、\passthrough{\lstinline![Bool]!}、\passthrough{\lstinline!Maybe Char!}；
\item
  态射是函数，如 \passthrough{\lstinline!Int -> Bool!}；
\item
  类型类是类型上的谓词，它筛选出一批类型，但不是范畴论意义上的对象。
\end{itemize}

一个常见误解是把 \passthrough{\lstinline!Monad!} 当作某种范畴。单子是 Hask 上的一个\textbf{自函子}（endofunctor）加两个
自然变换，不是对象的集合。

\subsection{以为恒等必须显式写出}\label{ux4ee5ux4e3aux6052ux7b49ux5fc5ux987bux663eux5f0fux5199ux51fa}

\passthrough{\lstinline!id!} 常常不出现，但恒等箭头始终存在，它由对象确定。\passthrough{\lstinline!identity :: Obj -> Arrow!} 表达的就是
这一点：给定对象，恒等随即确定。

这也解释了 \passthrough{\lstinline!fmap id = id!} 为何是定律：函子必须把恒等映射到恒等，这是定义的组成部分。

\subsection{Hask 并不严格满足公理}\label{hask-ux5e76ux4e0dux4e25ux683cux6ee1ux8db3ux516cux7406}

\passthrough{\lstinline!undefined :: Int!} 与 \passthrough{\lstinline!seq!} 会使部分等式失效：\passthrough{\lstinline!id . undefined!} 与 \passthrough{\lstinline!undefined!} 在 \passthrough{\lstinline!seq!} 下可
区分。严格地说，Hask 不满足单位律；结合律仍然成立，因为复合的两侧都化简为同一个 λ（见 6.3.2）。

本书在使用“Hask 是范畴”时，指的是忽略 \(\bot\) 的近似意义。第 6 章专门讨论这一问题。

\section{练习}\label{ux7ec3ux4e60}

\begin{enumerate}
\def\labelenumi{\arabic{enumi}.}
\tightlist
\item
  增加对象 \(D\) 与箭头 \(C \to D\)，补全 \passthrough{\lstinline!table!}，使 \passthrough{\lstinline!assocHolds!} 仍为 \passthrough{\lstinline!True!}。
\item
  改为菱形结构：\(A \to B\)、\(A \to C\)、\(B \to D\)、\(C \to D\)。\(A\) 到 \(D\) 有两条路径，但
  两条路径是否给出两个箭头，取决于复合的定义：若复合把两条路径算作同一个箭头，这仍是
  偏序范畴（例：集合包含关系，\(A \subseteq B \subseteq D\) 与 \(A \subseteq C \subseteq D\)
  给出同一个箭头 \(A \to D\)）；若复合保留路径，则 \(\mathrm{Hom}(A, D)\) 有两个元素。用代码
  表示后一种结构，并判断结合律是否仍成立。
\item
  补全恒等唯一性的证明，标出每一步所用的公理。
\item
  证明：若 \(F\) 是函子、\(f\) 是同构，则 \(Ff\) 也是同构。
\item
  用 \passthrough{\lstinline!Control.Category!} 的 \passthrough{\lstinline!Category!} 类构造实例：态射取 \passthrough{\lstinline!a -> Maybe b!}，即 Kleisli 箭头，
  \passthrough{\lstinline!id!} 取 \passthrough{\lstinline!return!}，复合取 \passthrough{\lstinline!(>=>)!}，并验证三条定律。这一构造即第 9 章的 Kleisli 范畴。
\end{enumerate}

\section*{参考文献}
\addcontentsline{toc}{section}{参考文献}\label{ux53c2ux8003ux6587ux732e}

本章的写法参考了以下著作，标注的是写作时实际核对的位置：

\begin{itemize}
\tightlist
\item
  \textbf{Awodey, Steve.} \emph{Category Theory}, 2nd ed. Oxford University Press, 2006.
  §1.3（范畴的定义）、§1.4（例子：偏序集、幺半群等）、§1.5（同构）。本章定义 5.1--5.3 的
  表述基本沿用这几节。
\item
  \textbf{Mac Lane, Saunders.} \emph{Categories for the Working Mathematician}, 2nd ed., Graduate Texts
  in Mathematics 5, Springer, 1998. Chapter I。本章未采用其记号，但“范畴 → 函子 → 自然变换”
  的排列顺序与之一致。
\item
  \textbf{\passthrough{\lstinline!Control.Category!}}（\passthrough{\lstinline!base!} 包）。\passthrough{\lstinline!id!} 与 \passthrough{\lstinline!(.)!} 的类，5.1 节的三条等式即其 API。
\end{itemize}

后续阅读顺序：Awodey §2（抽象结构：积、余积、初始与终止对象），再读 §6（指数对象、笛卡尔
闭范畴），分别对应本书第 7、8 章。
